\section{Cayley tables and complete mappings}\label{sec:complete-mappings}

The most natural Latin squares come from groups.

\begin{definition}\label{def:cayley-table}
The \term{Cayley table} of a finite group $G$ has its rows and columns indexed by $G$ and the entry $ab$ in cell $(a,b)$. It is a Latin square. As a colouring of a complete bipartite graph, it is the graph $K_{G,G}$ whose two sides are copies of $G$, with the edge $(x,y)$ coloured $xy$.
\end{definition}

The square of ranks in Ozanam's puzzle is a Cayley table. Renaming the ranks A, Q, J, K as $00,01,10,11$ and labelling the rows $00,11,01,10$ and the columns $00,01,10,11$, every cell becomes the sum of its row label and its column label in $\Z_2\times\Z_2$, and no cell moves (Figure~\ref{fig:ranks-cayley}).

\begin{figure}[ht]
\centering
\begin{tikzpicture}[every node/.style={font=\normalsize},
      cell/.style={draw, minimum size=7mm, inner sep=0pt, outer sep=0pt}]
    \def\cA{RainbowRed}\def\cK{RainbowBlue}\def\cQ{RainbowGreen}\def\cJ{black}%
    \def\gA{00}\def\gQ{01}\def\gJ{10}\def\gK{11}%
    \foreach \i/\row in {1/{A,Q,J,K}, 2/{K,J,Q,A}, 3/{Q,A,K,J}, 4/{J,K,A,Q}}{%
      \foreach \r [count=\j] in \row
        \node[cell] at (0.7*\j, -0.7*\i) {\textcolor{\csname c\r\endcsname}{\r}};}
    \node[font=\footnotesize] at (1.75,-3.45) {the ranks};
    \node[align=center, font=\small] at (4.6,-1.75)
      {$\xrightarrow{\ \text{rename}\ }$\\[4pt]
       \textcolor{RainbowRed}{A}$\,\mapsto00$\quad \textcolor{RainbowGreen}{Q}$\,\mapsto01$\\
       J$\,\mapsto10$\quad \textcolor{RainbowBlue}{K}$\,\mapsto11$};
    \begin{scope}[shift={(6.3,0)}]
      \foreach \i/\x/\row in {1/00/{A,Q,J,K}, 2/11/{K,J,Q,A},
                              3/01/{Q,A,K,J}, 4/10/{J,K,A,Q}}{%
        \node[font=\footnotesize, text=CiteColor] at (-0.02, -0.7*\i) {$\x$};
        \foreach \r [count=\j] in \row
          \node[cell] at (0.7*\j, -0.7*\i)
            {\textcolor{\csname c\r\endcsname}{$\csname g\r\endcsname$}};}
      \foreach \y [count=\j] in {00,01,10,11}
        \node[font=\footnotesize, text=CiteColor] at (0.7*\j, -0.2) {$\y$};
      \node[font=\footnotesize] at (-0.02,-0.2) {$+$};
      \node[font=\footnotesize] at (1.75,-3.45)
        {the addition table of $\mathbb Z_2\times\mathbb Z_2$};
    \end{scope}
  \end{tikzpicture}

\caption{Renaming the ranks and labelling the rows and columns by elements of $\Z_2\times\Z_2$, every cell becomes row $+$ column.}
\label{fig:ranks-cayley}
\end{figure}

So we may ask which groups $G$ have a Cayley table that decomposes into transversals. For a group, one transversal is enough.

\begin{fact}\label{fact:translates}
Let $T=\{(x,\phi(x)):x\in G\}$ be a transversal of the Cayley table of $G$. For each $g\in G$ the set
\[
T_g=\{(x,\phi(x)g):x\in G\}
\]
is again a transversal, and the $|G|$ sets $T_g$ partition the table.
\end{fact}

The entries of $T_g$ are those of $T$ multiplied on the right by $g$; in Ozanam's puzzle the four suits are exactly the four translates of the spade transversal (Figure~\ref{fig:translates}). For a group, therefore, Euler's question asks for a single rainbow perfect matching.

\begin{figure}[ht]
\centering
\begin{tikzpicture}[every node/.style={font=\normalsize},
      cell/.style={draw, minimum size=8mm, inner sep=0pt, outer sep=0pt}]
    \colorlet{suitS}{RainbowBlue!35}\colorlet{suitH}{RainbowRed!35}%
    \colorlet{suitD}{RainbowAmber!45}\colorlet{suitC}{RainbowGreen!40}%
    \foreach \i/\x/\row in {1/00/{00/S,01/C,10/H,11/D}, 2/11/{11/H,10/D,01/S,00/C},
                            3/01/{01/D,00/H,11/C,10/S}, 4/10/{10/C,11/S,00/D,01/H}}{%
      \node[font=\footnotesize, text=CiteColor] at (0.0, -0.8*\i) {$\x$};
      \foreach \v/\u [count=\j] in \row
        \node[cell, fill=suit\u] at (0.8*\j, -0.8*\i) {$\v$};}
    \foreach \y [count=\j] in {00,01,10,11}
      \node[font=\footnotesize, text=CiteColor] at (0.8*\j, -0.2) {$\y$};
    \node[font=\footnotesize] at (0.0,-0.2) {$+$};
    \foreach \u/\sym/\g/\y in {S/{$\spadesuit$}/{$T$}/-0.8,
        H/{\textcolor{RainbowRed}{$\heartsuit$}}/{$T+10$}/-1.6,
        D/{\textcolor{RainbowRed}{$\diamondsuit$}}/{$T+11$}/-2.4,
        C/{$\clubsuit$}/{$T+01$}/-3.2}{%
      \node[cell, fill=suit\u, minimum size=5.5mm] at (4.3,\y) {};
      \node[anchor=west] at (4.65,\y) {\sym\ \ \g};}
  \end{tikzpicture}

\caption{The four suits are the translates of one transversal $T$, obtained by shifting every column label by the same element.}
\label{fig:translates}
\end{figure}

\topic{Mann's complete mappings}

By 1942 every known pair of orthogonal Latin squares, whether in the statistical tables of Fisher and Yates~\cite{fy} or constructed from Galois fields by Bose~\cite{bose38} and Stevens~\cite{ste39}, was based on a group. Mann~\cite{man42} asked which groups can be used, and introduced the following notion.

\begin{definition}[Mann~\cite{man42}]\label{def:complete-mapping}
A bijection $\theta\colon G\to G$ is a \term{complete mapping} if $x\mapsto x\,\theta(x)$ is also a bijection.
\end{definition}

In the Cayley table, a complete mapping chooses in row $x$ the cell in column $\theta(x)$; these $n$ cells form a transversal, and every transversal arises in this way.

\begin{theorem}[Mann~\cite{man42}; Paige~\cite{pai51}]\label{thm:mann}
The Cayley table of $G$ has an orthogonal mate if and only if $G$ has a complete mapping $\theta$. One mate is the Latin square $S(i,k)=i\,k\,\theta(k)$.
\end{theorem}

The translation to rainbow matchings is immediate.

\begin{proposition}\label{prop:cm}
A bijection $\theta\colon G\to G$ is a complete mapping if and only if the perfect matching $\{(x,\theta(x)):x\in G\}$ of $K_{G,G}$ is rainbow. Every rainbow perfect matching of $K_{G,G}$ arises in this way from exactly one $\theta$.
\end{proposition}

\begin{proof}
A perfect matching of $K_{G,G}$ is the graph of a bijection $\theta$. Its colours are the elements $x\theta(x)$, and these are distinct exactly when $x\mapsto x\theta(x)$ is injective, that is, a bijection.
\end{proof}

\begin{corollary}\label{cor:cm}
The following are equivalent: $G$ has a complete mapping; $K_{G,G}$ has a rainbow perfect matching; the edges of $K_{G,G}$ split into $|G|$ rainbow perfect matchings.
\end{corollary}

\begin{proof}
Combine Proposition~\ref{prop:cm} with Fact~\ref{fact:translates} and Proposition~\ref{prop:mate}.
\end{proof}

\topic{The first results}

The first results on complete mappings appeared between 1942 and 1951.
\begin{itemize}
\item Mann~\cite{man42} showed in 1942 that \emph{no group of order $4n+2$ has a complete mapping}, so groups can never refute Euler's conjecture. In the same paper he obtained new orthogonal pairs of order $12$, from $\Z_2\times\Z_2\times\Z_3$ and from $A_4$, beyond the reach of the Galois-field method.
\item Paige~\cite{pai47} showed in 1947 that an abelian group has a complete mapping if and only if $\sum_{x\in G}x=0$.
\item Bateman~\cite{bat50} showed in 1950 that every infinite group has a complete mapping.
\item Paige~\cite{pai51} showed in 1951 that every group of odd order has one, namely $\theta(x)=x$.
\end{itemize}
Paige also found a necessary condition for general groups.

\begin{theorem}[Paige~\cite{pai51}]\label{thm:paige}
If $G$ has a complete mapping, then the product of all elements of $G$, taken in any order, lies in the commutator subgroup $G'$.
\end{theorem}

Paige conjectured the converse. He also conjectured that no symmetric group $S_n$ has a complete mapping.

\begin{example}\label{ex:order-four}
The two groups of order $4$ already behave differently (Figure~\ref{fig:z2z2-z4}). In $\Z_2\times\Z_2$ we have $00+01+10+11=00$, so Paige's condition holds, and the Cayley table has a transversal. In $\Z_4$ we have $0+1+2+3=2\neq0$, so Paige's condition fails and the Cayley table has no transversal. The two tables are Latin squares of the same order; what separates them is the structure of the group.
\end{example}

\begin{figure}[ht]
\centering
\begin{tikzpicture}[every node/.style={font=\normalsize},
      cell/.style={draw, minimum size=7.5mm, inner sep=0pt, outer sep=0pt}]
    \def\cA{RainbowRed}\def\cK{RainbowBlue}\def\cQ{RainbowGreen}\def\cJ{black}%
    \def\gA{00}\def\gQ{01}\def\gJ{10}\def\gK{11}%
    \foreach \i/\x/\row in {1/00/{A,Q,J,K}, 2/11/{K,J,Q,A},
                            3/01/{Q,A,K,J}, 4/10/{J,K,A,Q}}{%
      \node[font=\footnotesize, text=CiteColor] at (0.0, -0.75*\i) {$\x$};
      \foreach \r [count=\j] in \row
        \node[cell] at (0.75*\j, -0.75*\i)
          {\textcolor{\csname c\r\endcsname}{$\csname g\r\endcsname$}};}
    \foreach \y [count=\j] in {00,01,10,11}
      \node[font=\footnotesize, text=CiteColor] at (0.75*\j, -0.2) {$\y$};
    \node[font=\footnotesize] at (0.0,-0.2) {$+$};
    \foreach \i/\j in {1/1, 2/3, 3/4, 4/2}
      \node[circled, minimum size=6.6mm] at (0.75*\j, -0.75*\i) {};
    \node at (1.875,-3.85) {$\mathbb Z_2\times\mathbb Z_2$};
    \begin{scope}[shift={(5.2,0)}]
      \foreach \i in {0,...,3}{%
        \node[font=\footnotesize, text=CiteColor] at (0.0, -0.75*\i-0.75) {$\i$};
        \node[font=\footnotesize, text=CiteColor] at (0.75*\i+0.75, -0.2) {$\i$};
        \foreach \j in {0,...,3}{%
          \pgfmathtruncatemacro{\v}{mod(\i+\j,4)}%
          \node[cell] at (0.75*\j+0.75, -0.75*\i-0.75) {$\v$};}}
      \node[font=\footnotesize] at (0.0,-0.2) {$+$};
      \node at (1.875,-3.85) {$\mathbb Z_4$};
    \end{scope}
  \end{tikzpicture}

\caption{The Cayley tables of $\Z_2\times\Z_2$ and $\Z_4$. The first has a transversal (circled); the second has none.}
\label{fig:z2z2-z4}
\end{figure}

\topic{Hall and Paige, 1955}

Hall and Paige~\cite{hp55} translated Paige's condition into a statement about Sylow subgroups.

\begin{fact}[Hall--Paige~\cite{hp55}]\label{fact:sylow}
The product of all elements of $G$ lies in $G'$ if and only if the Sylow $2$-subgroups of $G$ are trivial or non-cyclic.
\end{fact}

For $\Z_4$ the Sylow $2$-subgroup is $\Z_4$ itself, which is cyclic; for $\Z_2\times\Z_2$ it is not. In the same paper Hall and Paige proved that every \emph{solvable} group with this property has a complete mapping, and that $S_n$ for $n>3$ and $A_n$ have complete mappings, so Paige's conjecture on symmetric groups was false. They conjectured the converse of Theorem~\ref{thm:paige} in the following form.

\begin{conjecture}[Hall--Paige~\cite{hp55}]\label{conj:hp}
A finite group $G$ has a complete mapping if and only if its Sylow $2$-subgroups are trivial or non-cyclic.
\end{conjecture}

We say that $G$ satisfies the \term{Hall--Paige condition} if its Sylow $2$-subgroups are trivial or non-cyclic.
